\SourceLeaf{285}%
\(GF\) in \(G\) bezeichnet), und bezeichnet man ebenso die (immer negativen) Werte von \(\tfrac{206265''\mathrm{d}m}{2m\mathrm{d}S}\) in denselben Punkten mit \(-k^0,-k'\), so erhält man
\[
\begin{aligned}
206265''l^0&=-2k^0\sin V^0\\
206265''l'&=+2k'\sin V'
\end{aligned}
\]

Rechnet man die im vorhergehenden Artikel angegebenen Ausdrücke für \(\psi^0,\psi'\) in Sekunden um und lässt die Faktoren \(\BareRoot[6]{\tfrac{m^0}{m'}},\BareRoot[6]{\tfrac{m'}{m^0}}\) weg, die hier nur unmerklich von eins abweichen, so erhält man
\[
\begin{aligned}
\psi^0&=-\tfrac{1}{3}h(2k^0\sin V^0-k'\sin V')\\
\psi'&=-\tfrac{1}{3}h(2k'\sin V'-k^0\sin V^0)
\end{aligned}
\]

Die dieser Abhandlung beigefügte Tafel gibt in der letzten Spalte unter der Überschrift \(k\) die Werte von \(k^0,k'\) für die entsprechenden Werte von \(S\) an, die in der ersten Spalte unter der Überschrift \(Q+q\) zu suchen sind. Da \(k\) stets positiv ist und \(\sin V^0,\sin V'\) immer entgegengesetzte Vorzeichen haben, ist \(\psi^0\) negativ und \(\psi'\) positiv, wenn \(G\) westlich von \(F\) liegt, und umgekehrt. Bei der Rechnung ist zu beachten, dass \(h\) in diesen Formeln in Teilen des Radius angegeben ist. Der in einem beliebigen Längenmaß angegebene Abstand der Punkte \(F,G\) ist deshalb zuvor durch den im gleichen Maß ausgedrückten Wert von \(A\) zu teilen.

\begingroup\rightskip=0pt plus 1em
Bei unserer konformen Abbildung der Ellipsoidfläche auf die Kugelfläche wird ein Meridian der ersten Fläche wieder durch einen Meridian der zweiten dargestellt. Somit hat jedes Element von \(L\) dieselbe Neigung zum Meridian wie das entsprechende Element von \(M\). Folglich sind die Azimute der geodätischen Linie in ihren beiden Endpunkten \(V^0+\psi^0\) beziehungsweise \(V'+\psi'\). Sind umgekehrt diese Azimute gegeben, so erhält man die auf die Kugelfläche reduzierten Azimute, indem man \(-\psi^0,-\psi'\) anbringt. Bei der Berechnung dieser stets fast verschwindenden Reduktionen macht es offensichtlich keinen Unterschied, wenn man in den obigen Formeln anstelle von \(V^0,V'\) die Azimute auf dem Ellipsoid verwendet.\par\endgroup

\Article{14}
Um nach den angegebenen Vorschriften die Reduktionen der Richtungen für die Übertragung vom Ellipsoid auf die Kugel oder umgekehrt berechnen zu
